\subsection{The full source and its relative completions}

We write
\[
 \mathcal B=\widehat{C_2*C_3}
   =\langle s,r\mid s^2=r^3=1\rangle^\wedge,\qquad
 \tau=sr,\qquad F=\ker(\mathcal B\to S_3).
\]
All profinite commutator subgroups are closed. For an open normal
subgroup \(U\leq\Gamma\), put
\[
 R_{\rm nil}(U)=
 \bigcap_{\substack{\phi:U\to L\ {\rm continuous}\\
                         L\ {\rm finite\ nilpotent}}}\ker\phi,
 \qquad \Gamma_{U,\rm nil}=\Gamma/R_{\rm nil}(U).
\]
This is a relative completion with its specified map to \(\Gamma/U\);
it is not the pronilpotent completion of \(\Gamma\).
A positive peripheral automorphism sends each named peripheral
generator to a conjugate of itself. We retain genuine actions and
three actual conjugators until explicitly taking outer classes.
We read LM7\_PUSH D.50, E.58 and the attached referee report.
The following reconstruction checks the proof, rather than relying
on that report's verdict \citep{cLM7Push,cLM7}. The auxiliary positive-genus formula in
LM7\_PUSH D.50 is unnecessary and is not asserted here.

\begin{statement}{C.1}{Finite central refinements retain generation and conjugators}{Proved}
Let \(\Lambda\) be the free profinite group on \(x_1,x_2\), with
\(x_3=(x_1x_2)^{-1}\), or \(\mathcal B\), with
\((x_1,x_2,x_3)=(s,r,\tau^{-1})\).
Given finite continuous quotients
\[
 q_E:\Lambda\twoheadrightarrow E,\qquad
 1\to A\to E\xrightarrow{\pi}D\to1,\qquad q_0=\pi q_E,
\]
there are a finite central extension \(p:D^\dagger\twoheadrightarrow D\)
and an actual epimorphism \(q^\dagger:\Lambda\twoheadrightarrow D^\dagger\)
above \(q_0\) with this property: every epimorphism
\(q_1^\dagger:\Lambda\twoheadrightarrow D^\dagger\) satisfying
\[
 q_1^\dagger(x_i)=b_i^\dagger q^\dagger(x_i)(b_i^\dagger)^{-1}
\]
has an epimorphism lift \(f:\Lambda\twoheadrightarrow E\) of
\(pq_1^\dagger\), with peripheral conjugators whose images in \(D\)
are exactly \(p(b_i^\dagger)\).
\end{statement}
\begin{proof}
Write \(X_i=q_E(x_i)\), and let \(T_i\) be its \(A\)-conjugacy class.
Retain every \(J\leq E\) onto \(D\) with all \(T_i\cap J\ne\varnothing\).
Put \(M=J\cap A\), \(e=\exp J\). For every \(J\)-invariant irreducible
\(\rho:M\to{\rm GL}_d(\mathbb C)\), set
\[
 \widetilde J_\rho=
 \left\{(g,B):
 \begin{array}{l}
 g\in J,\ B\rho(m)B^{-1}=\rho(gmg^{-1})\ (m\in M),\\
 \det B\in\mu_e
 \end{array}\right\}.
 \tag{C.1a}
\]
Schur's lemma makes each fibre have \(de\) elements, with central
scalar kernel \(\mu_{de}\). The graph \(m\mapsto(m,\rho(m))\) is
normal, lies in this finite group, and meets that kernel trivially.
Consequently
\[
 Q_\rho=\widetilde J_\rho/\operatorname{graph}(\rho)\to D
 \tag{C.1b}
\]
is a finite central extension. The map
\(\widetilde J_\rho\to J\times_DQ_\rho\) sending
\((g,B)\) to \((g,[(g,B)])\) is an isomorphism: a graph element
corrects the first coordinate of any chosen representative.
Its matrix coordinate gives a representation \(\Psi_\rho\)
of this actual pullback.

There is a homomorphism \(R_\rho^0:\Lambda\to Q_\rho\) above \(q_0\).
For the free source choose arbitrary generator lifts. In the modular
case choose \(a_s\in T_1\cap J,a_r\in T_2\cap J\).
Their implementing matrices have scalar square and cube, respectively.
Rescale them to have square and cube one. The determinants then lie
in \(\mu_2,\mu_3\subset\mu_e\); if the element is one, take the
identity matrix instead. Thus these choices respect the actual
torsion relations. Define \(D^\dagger\) to be the joint image of
\(q_0\) and these finitely many maps. Its kernel over \(D\) is central.

We prove the count needed to preserve surjectivity. Let
\(\chi:\widehat F_2\twoheadrightarrow D\), and let
\(U_i\subset J_{\chi(x_i)}\) be invariant under \(M\)-conjugacy.
Choose actual lifts \(R_\rho:\widehat F_2\to Q_\rho\) of \(\chi\).
The number of homomorphisms to \(J\) with these values is
\[
 \frac1{|M|}
 \sum_{\substack{\rho\in{\rm Irr}(M)\\J\ {\rm fixes}\ [\rho]}}
 \frac1{\dim\rho}
 \prod_{i=1}^3\left(
  \sum_{a\in U_i}{\rm tr}\,\Psi_\rho(a,R_\rho(x_i))\right).
 \tag{C.1c}
\]
Indeed \(B_i=\sum_{a\in U_i}a\in\mathbb C[J]\) commutes with
\(\mathbb C[M]\). If the quotient value of \(x_i\) moves a central
primitive idempotent \(e_\rho\), then
\(e_\rho B_i=e_\rho B_ie_\rho=0\). The first two quotient values
generate \(D\), so only invariant irreducibles survive in the product.
For those irreducibles each lifted matrix sum is scalar.
The product of the three second coordinates is one in \(Q_\rho\),
so its matrix is literally \(\rho(a_1a_2a_3)\).
The regular-character identity
\[
 {\bf1}_{\{1\}}(m)=|M|^{-1}
       \sum_\rho(\dim\rho){\rm tr}\rho(m)
\]
gives (C.1c), including its exponent \((\dim\rho)^{-1}\).
There is no omitted projective scalar. In the modular case the first
two peripheral classes impose \(s^2=r^3=1\), so the same count
counts maps from the actual modular source.

Put \(h_i=p(b_i^\dagger)\), choose \(H_i\in E\) over \(h_i\),
and set \(T_i^1=H_iT_iH_i^{-1}\).
For \(J\) onto \(D\), choose \(j_i\in J\) over \(h_i\).
As \(H_ij_i^{-1}\in A\), we have exactly
\[
 T_i^1\cap J=j_i(T_i\cap J)j_i^{-1}.
 \tag{C.1d}
\]
Thus omitted subgroups have count zero on both sides.
The pair \((j_i,\overline b_i^\dagger)\in J\times_DQ_\rho\)
conjugates the old lifted boundary matrices to the new ones.
Every trace, and hence the homomorphism count into every \(J\),
is unchanged. Assign count zero to subgroups not onto \(D\).
Inversion in the finite subgroup lattice gives
\[
 \#\{\hbox{epimorphisms in the fibre}\}
   =\sum_{J\le E}\mu(J,E)
          \#\{\hbox{homomorphisms in the fibre into }J\}.
 \tag{C.1e}
\]
The epimorphism counts agree, and the old count contains \(q_E\).
Each element of \(T_i^1\) is conjugate to \(X_i\) by an element
of \(H_iA\), giving exactly the prescribed conjugator image.
\end{proof}

\begin{statement}{C.2}{Whole relative pronilpotent lifting}{Proved}
For either source \(\Gamma\) in C.1 and every open normal \(U\),
each continuous positive peripheral automorphism of
\(\Gamma_{U,\rm nil}\), with three specified conjugators there,
lifts to a continuous automorphism of \(\Gamma\), retaining
those conjugators literally. If one specified conjugator is one,
the corresponding peripheral can be fixed literally.
\end{statement}
\begin{proof}
For any finite continuous quotient \(q\) of a profinite group \(V\),
\[
 R_{\rm nil}(V)=\bigcap_n\gamma_nV,\qquad
 q(R_{\rm nil}(V))=\gamma_\infty q(V).
 \tag{C.2a}
\]
For the first equality, bounded-class nilpotent profinite quotients
separate elements; for the second use compactness of the descending
fibres in the closed \(\gamma_nV\).
Work in \(\Gamma^3\), the compact space of conjugators.
Impose the peripheral product relation, generation of every finite
quotient, and all finite coordinates of the specified conjugators
in \(\Gamma_{U,\rm nil}\). Any finite list of these closed conditions
is contained in one finite quotient \(E_0\), enlarged to retain
\(\Gamma/U\). Put
\[
 U_0=\operatorname{im}(U\to E_0),\quad
 A_0=\gamma_\infty U_0,\quad D_0=E_0/A_0.
 \tag{C.2b}
\]
Every old relative constraint kills \(A_0\).
The image of \(U\) in the C.1 refinement \(D_0^\dagger\)
is a central extension of the nilpotent \(U_0/A_0\), hence
nilpotent. Thus its actual source map factors through
\(\Gamma_{U,\rm nil}\). Precompose with the given automorphism
\(\alpha\). C.1 gives an epimorphism to \(E_0\) and conjugators
with every old coordinate. Lifting those three elements to
\(\Gamma\) proves the finite intersection property.

The resulting conjugators define an endomorphism of the full source.
For \(\mathcal B\), conjugates of \(s,r\) satisfy their torsion
relations literally, so the universal property applies.
Its compact image generates every finite quotient and is therefore
the whole group. An onto endomorphism of a finitely generated
profinite group is injective: for each index preimage permutes the
finite set of open subgroups. Periodicity of each such subgroup
forces the endomorphism's kernel into it; their intersection is one.
This is the automorphism step, not merely an endomorphism construction.
The retained map agrees with \(\alpha\) on generators.
If a specified conjugator is one, its lift lies in
\(R_{\rm nil}(U)\); inner conjugation by its inverse fixes that
peripheral literally and preserves all retained coordinates.
This reconstructs LM7\_PUSH D.50 and the relative part of E.58.
\end{proof}

\subsection{An arithmetic input and the fixed-fibre comparison}

\begin{statement}{C.3}{A certified cubic CM commutator}{Proved}
Put \(k=\mathbb Q(\zeta)\), \(\zeta^2+\zeta+1=0\),
\(E_3:y^2=x^3+1\), \(E_4:y^2=x^3-x\).
The image of \(G_k\) on \(T_2E_3\), viewed as a free rank-one
\(\mathcal O=\mathbb Z_2[\zeta]\)-module, is \(1+2\mathcal O\).
There is \(h\in[G_{\mathbb Q},G_{\mathbb Q}]\) with identity
genuine tangential pro-two tripod action and CM actions
\((-I,[i])\) on \(E_3[4]\) and projective \(E_4[3]\).
Its normalized congruence action on
\(Q={\rm SL}_2(\mathbb Z/12)/\{\pm I\}\) is the identity.
\end{statement}
\begin{proof}
The automorphism \((x,y)\mapsto(\zeta x,y)\) is defined over \(k\);
complex uniformization identifies the integral Tate module with
a rank-one \(\mathcal O\)-module. All two-torsion is rational over
\(k\), so the image is contained in \(U_2=1+2\mathcal O\).
Exact point counting gives \(a_7=-4,a_{13}=2\).
The Frobenius trace and determinant formula
\citep[V.7.5 and its preceding determinant formula]{cMilne2021}
applies at these split good primes, different from 2.
The respective CM roots and their conjugates are
\[
 \{-1+2\zeta,-3-2\zeta\},\qquad
 \{3+4\zeta,-1-4\zeta\}.
 \tag{C.3a}
\]
Conjugation supplies both values since the curve is over \(\mathbb Q\).
The program \path{c_cubic_cm_image.py} independently counts
the points and verifies that these roots generate all 16 units in
\(U_2/(1+8\mathcal O)\). Moreover
\(1+8\mathcal O\subset U_2^2\), using
\(\exp(\tfrac12\log u)\in1+4\mathcal O\).
Thus the closed image surjects to \(U_2/U_2^2\).
A proper closed subgroup of an abelian pro-two group is detected
by a quotient of order two, so the image is all \(U_2\).

The cyclic cubic field \(k(\sqrt[3]4)\) is disjoint from that
pro-two division field. Choose \(g\in G_k\) acting by
\(a=1-2\zeta^2=3+2\zeta\) on \(T_2E_3\) and fixing \(\sqrt[3]4\).
Take \(\tau_3\in I_3\) nontrivial on \(k\), and put
\(h=g\tau_3g^{-1}\tau_3^{-1}\).
The genuine pro-two action based at \(\overrightarrow{01}\) is
unramified outside 2 \citep[\S8.5]{cHainMatsumoto2003}, so
\(\tau_3\), and then this commutator, act identically.
CM semilinearity gives
\[
 \rho_{E_3,2}(h)=a/\bar a=(-3-8\zeta^2)/7\equiv-1\pmod4.
 \tag{C.3b}
\]

We check LM7 D.25's finite field comparison.
For \(r_0=(3/4)^{1/4}\), halves of the two-torsion point
with \(x=e\in\{-1,-\zeta,-\zeta^2\}\) have
\[
 x=e(1\pm\sqrt3),\qquad y^2=-9\pm6\sqrt3.
 \tag{C.3c}
\]
A square root for the plus sign is
\(\sqrt3(\sqrt3-1)r_0\).
These formulas and their inverses give
\(\mathbb Q(E_3[4])=\mathbb Q(i,r_0)\).
For \(E_4\), the three-division polynomial is
\(3x^4-6x^2-1\).
Choose roots \(a_0,b_0\) with
\(a_0^2=1+2/\sqrt3\), \(b_0^2=1-2/\sqrt3\),
\(a_0b_0=i/\sqrt3\).
Then
\[
 w=a_0+ib_0,\quad w^4=4/3,\quad
 a_0=(w+2/w)/2,\quad b_0=(w-2/w)/(2i).
 \tag{C.3d}
\]
Thus \(\mathbb Q(x(E_4[3]))=\mathbb Q(i,r_0)\).
The commutator fixes every abelian extension.
By (C.3b)--(C.3c) it sends \(r_0\) to \(-r_0\); by
(C.3d) it negates the nonzero three-torsion \(x\)-coordinates.
This is the projective action of multiplication by \(i\).

For the Kummer cocycle \(\kappa_2\), its two-primary part
is unramified at 3, so
\(\kappa_{2,2}(\tau_3)=0,\chi_2(\tau_3)=1\).
The commutator cocycle formula gives \(\kappa_{2,2}(h)=0\).
Modulo 3, \(\chi_3(g)=1\), and fixing \(\sqrt[3]4\) gives
\(\kappa_{2,3}(g)=0\); hence \(\kappa_{2,3}(h)=0\pmod3\).
Marin's normalized congruence shear has parameter
\(-8\kappa_2\) and diagonal \(\chi\)
\citep[\S4, Lemmas 4.1--4.3 and the arithmetic specialization]{cMarin2008}.
Its reduction modulo 12 is the identity, as claimed.
\end{proof}

\begin{statement}{C.4}{The compatible conjugators of LM7 D.46}{Proved}
Let \(N=\ker(\mathcal B\to Q)\) and
\(\mathcal P_p=\mathcal B/O^p(N)\).
For the normalized positive Galois action of C.3,
\[
 \mathcal U=\mathcal B/R_{\rm nil}(N)
       \simeq\mathop{\prod}_{p,Q}\mathcal P_p,\qquad
 \alpha^\dagger=(\alpha_{h,2},1,1,\ldots)
 \tag{C.4a}
\]
is an actual continuous automorphism with three compatible
conjugators in \(\mathcal U\).
It satisfies C.2 for \(U=N\), and its normalized full lifts
are multiplier-one gentle elements.
Their CM pair, with the fixed-fibre and global-sign
conventions proved in C.5, is \((I,[i])\).
\end{statement}
\begin{proof}
The map \(N/R_{\rm nil}(N)\to\prod_pN^{(p)}\) is an
isomorphism: distinct prime quotients have no nontrivial
common finite quotient, and compactness gives surjectivity
onto their whole product. To obtain (C.4a), lift a tuple's
common \(Q\)-value to \(\mathcal B\), then correct in this
kernel product. All component actions induce the identity
on \(Q\), so the product and its inverse are continuous.

Take actual \(h\)-conjugators \(A_h,\beta_h\) for \(r,s\);
the cusp conjugator is one. Their \(Q\)-values \(a_Q,b_Q\)
centralize \(q(r),q(s)\).
In a surjection with pro-\(p\) kernel, a finite-order
prime-to-\(p\) element's centralizer maps onto the centralizer
of its image. Indeed its cyclic subgroup and the conjugate
by a lift are complements in the same extension.
Coprime complement conjugacy at each finite level, then
compactness, supplies a correcting element in the kernel.
Thus choose \(c_2\in C_{\mathcal P_2}(r)\) over \(a_Q^{-1}\),
and \(b_p\in C_{\mathcal P_p}(s)\) over \(b_Q\) for every odd \(p\).
The actual tuples
\[
 A^\dagger=(A_{h,2}c_2,1,1,\ldots),\quad
 \beta^\dagger=(\beta_{h,2},b_3,b_5,\ldots),\quad V^\dagger=1
 \tag{C.4b}
\]
have common \(Q\)-values \(1,b_Q,1\) and implement the
peripheral relations. C.2 now applies literally.
After cusp normalization \(f=s\Phi(s)\in F\).
It is trivial in every pure pro-\(p\) quotient, by C.3 at
two and the identity component at odd primes.
In particular its abelianization is trivial, so \(f\in F'\).
The relations \((sf)^2=(f^{-1}r)^3=1\) are the normalized
gentle relations in \(x=\tau^{-2},y=sxs\).
The map is an automorphism by C.2.
The last assertion follows from the fixed-fibre geometric
comparison proved next.
\end{proof}

\begin{statement}{C.5}{Arithmetic frames and the CM pair of the full lift}{Proved}
Use the positive tangential base point at \(t=0\), the
marked modular generators of C.1, and the corresponding
transported Tate bases. For \(g\in G_{\mathbb Q^{\rm ab}}\),
the normalized pure modular frames satisfy
\[
 [q(A_g)^{-1}V_g]=\overline\rho_{E_3}(g),\qquad
 [q(\beta_g)^{-1}V_g]=\overline\rho_{E_4}(g).
 \tag{C.5a}
\]
The bars divide by the \emph{single global} subgroup
\(\{I,-I\}\), not independently by a sign at every prime.
Consequently C.4 gives full multiplier-one gentle
elements with fixed-fibre CM pair \((I,[i])\).
In particular the assertion of LM7 D.19(ii), namely
\(\Delta\ne0\) on \(\widehat{\rm GT}_0\), holds by this
construction as well as by the independent E.58
construction reconstructed in C.18.
\end{statement}
\begin{proof}
We supply the based comparison rather than inferring it
from equality of division fields. Put
\[
 \mathcal X=\mathcal M_{1,1}^{\rm rig}
       \simeq[(\mathbb P^1-\{0,1,\infty\})/S_3].
 \tag{C.5b}
\]
Here rigidification removes the central automorphism
\([-1]\). To verify the equivalence over \(\mathbb Q\),
quotient an elliptic curve by \([-1]\), mark the image
of its identity, and order the other three branch points.
The ordered four-pointed genus-zero curve has a Legendre
parameter; changing the order gives precisely the
anharmonic \(S_3\)-action. Conversely its double cover
recovers the elliptic curve fppf locally, with exactly
the central ambiguity \(\{\pm1\}\). These inverse
constructions on objects and isomorphisms descend after
rigidification. The modular orbifold and its universal
elliptic curve are described in
\citet[Example 3.14, \S3.6 and \S8.4]{cHain2010}.
We use the rigidified stack in (C.5b), not a direct-product
description of the unrigidified stabilizers.

The universal Tate local system on \(\mathcal M_{1,1}\)
therefore descends, after dividing by this one global
sign, to a continuous representation
\[
 \widetilde q:\pi_1(\mathcal X,\vec{01})
    \longrightarrow
       {\rm GL}_2(\widehat{\mathbb Z})/\{\pm I\}.
 \tag{C.5c}
\]
Its geometric representation is the usual modular map
\(q:\mathcal B\to{\rm SL}_2(\widehat{\mathbb Z})/\{\pm I\}\).
Indeed the integral homology local system has monodromy
\({\rm SL}_2(\mathbb Z)\); modulo the central sign its
two elliptic inertia generators and cusp give
\({\rm PSL}_2(\mathbb Z)=C_2*C_3\).
Reduction modulo every integer and profinite comparison
give the stated continuous geometric map.

Let \(\sigma_b\) be the tangential arithmetic section.
Its image under (C.5c) is exactly the class of the
normalized cusp matrix \(V_g\). Here one must exclude
independent prime-dependent signs. The Legendre invariant
gives
\[
 q_{
m Tate}=t^2/256\,(1+t\mathbb Q[[t]]).
 \tag{C.5d}
\]
The Tate torsion representatives are
\(\zeta_n^a q_{
m Tate}^{b/n}\); their construction
uses the universal Laurent-series identities over
\(\mathbb Z((q_{
m Tate}))\), which remain identities
after this substitution
\citep[III, pp.~102--103, equations (23)--(26)]{cMilne2021}.
The tangential section fixes the compatible roots of
\(t\) and the roots with constant term one of the
displayed unit. Its matrix is consequently
\(\left(\begin{smallmatrix}\chi(g)&-8\kappa_2(g)\\0&1
\end{smallmatrix}\right)\), up to the single quadratic
twist from the Legendre curve to the split Tate curve.
That twist is one character with values \(\{\pm1\}\)
on every Tate module, so disappears in (C.5c).
This also proves the required compatibility with the
normalized congruence matrix, including its scalar.
For the present subgroup \(\chi=1\), changing the
compatible cusp splitting by an integral shear changes
the upper entry by a multiple of \(1-\chi\), hence
does not change it.

The rational objects \(E_3\) and \(E_4\) give arithmetic
sections of (C.5b). Transport them to \(b\) by paths
whose elliptic inertia generators are the specified
\(r\) and \(s\), and call them \(\sigma_3,\sigma_4\).
For \(g\in G_{\mathbb Q^{\rm ab}}\), these sections
centralize the respective inertia generators: all the
CM automorphisms are defined over
\(\mathbb Q(\sqrt{-3})\) or \(\mathbb Q(i)\), both fixed
by \(g\). Define actual geometric elements by
\[
 A_g=\sigma_b(g)\sigma_3(g)^{-1},\qquad
 \beta_g=\sigma_b(g)\sigma_4(g)^{-1}.
 \tag{C.5e}
\]
They conjugate \(r,s\) exactly as the normalized
Galois action does. They are pure: the images of both
sections modulo two are identity, because
\(E_3[2]\) is defined over \(\mathbb Q(\sqrt{-3})\),
\(E_4[2]\) is rational, and \(V_g\equiv I\pmod2\).
The torsion centralizer theorem gives
\[
 C_{\mathcal B}(r)\cap F=C_{\mathcal B}(s)\cap F=1.
\]
We use \citet[Theorem B']{cHerfortRibes1985} here.
Thus (C.5e) are exactly the prescribed pure frames,
not merely other frames in the same outer class.
Applying (C.5c) to
\(\sigma_3(g)=A_g^{-1}\sigma_b(g)\) and its quartic
counterpart proves (C.5a).

For clarity, the familiar Legendre specializations give
the same fixed objects modulo this sign. At \(t=1/2\),
\(X=4x-2,Y=8y\) gives \(Y^2=X^3-4X\), the
quadratic twist of \(E_4\) by 2. At
\(t^2-t+1=0\), put \(d=2t-1\); shifting the
\(x\)-coordinate gives
\(y^2=x^3+d/9\), and \(d/9=(-d/3)^3\).
This is a quadratic twist of \(E_3\) over
\(\mathbb Q(\sqrt{-3})\). Thus no cubic twist has
silently been replaced by a quadratic one.

The rational point \((0,1)\) of order three on \(E_3\)
selects the unique global sign of its arithmetic Tate
representation. With this convention C.3 gives
\(T_h|_{E_3[4]}=-I\) and \(W_h|_{E_4[3]}=[i]\).
For the C.4 frames, however,
\(q(A^\dagger)=1\), \(q(\beta^\dagger)=b_Q\),
and the retained congruence action is identity.
The unique normalized shear commuting with the two
congruence braid generators is identity modulo 12.
Hence their cubic coordinate is identity even with
the rational-three-torsion sign, and their quartic
coordinate modulo three is \(b_Q^{-1}=W_h=[i]\).
All these finite coordinates survive C.2's lift.
Their E.65 bits are \((0,1)\), so \(\Delta=1\),
as also follows from the direct parity proof in C.18.
The construction uses no assertion that a nongalois
automorphism of a compact elliptic cover comes from
a fixed arithmetic fibre.
\end{proof}

\subsection{The two lifts and their residuals}

\begin{statement}{C.6}{Parity and the pronilpotent-kernel comparison}{Proved}
Use the CM identifications and character construction
proved in C.5 and C.18, reconstructing
LM7\_PUSH E.58 and E.65 (R13--R14).
For an E.58 lift \(e\) and a C.4 lift \(d\),
\[
 \Delta(e)=\Delta(d)=1,\qquad \Delta(de^{-1})=0.
 \tag{C.6a}
\]
Both induce identity on \(\mathcal B/R_{\rm nil}(F)\).
Their quotient therefore belongs already to this action
kernel, and hence to the closed subgroup generated by that
kernel and the Galois image.
\end{statement}
\begin{proof}
The two CM bits are \((1,0)\) and \((0,1)\).
E.65 defines their sum as the restriction of its character
\(\Delta\), proving (C.6a).
The first lift retains identity on
\(\mathcal B/R_{\rm nil}(H)\), \(H=F'F^3\).
Restricting a finite nilpotent quotient of \(F\) to \(H\)
gives a nilpotent group, so
\(R_{\rm nil}(H)\subset R_{\rm nil}(F)\).
For the second lift, each map
\(\mathcal B\to\mathcal B/O^p(F)\) kills \(O^p(N)\),
since the image of \(N\) is pro-\(p\).
The induced normalized action is identity at every \(p\),
as in C.4. Their fibre product over \(S_3\) is
\(\mathcal B/R_{\rm nil}(F)\).
No assertion that \(de^{-1}\) is itself Galois enters
this proof.
\end{proof}

\begin{statement}{C.7}{The forced first homology images and Lead 6}{Open}
For \(G_j=\mathcal B/[N,N]N^j\), the following part of the
images is determined for \emph{every} choice of the full lifts:
\[
\begin{array}{c|c|c|c}
 &G_2&G_3&G_6=G_2\mathbin{\times}_QG_3\\ \hline
 e&1&\varepsilon_3\ne1&(1,\varepsilon_3)\\
 d&h_2\ne1&1&(h_2,1)\\
 de^{-1}&h_2\ne1&\varepsilon_3^{-1}\ne1&(h_2,\varepsilon_3^{-1})
\end{array}
\tag{C.7a}
\]
Here \(\varepsilon_3\) is the actual induced action of the
chosen E.58 lift, not an asserted unique automorphism of \(G_3\).
The complete matrix of this action and the equality with a
Galois action in task C.2 and Lead 6 remain undetermined.
\end{statement}

\begin{proof}[Proof of the asserted finite information]
The image of \(H=F'F^3\) in \(G_2\) is a two-group:
its image in \(Q\) is \(V_4\times Q_8\), and its kernel
is a two-group. Hence \(G_2\) factors through the E.58
retained relative source, on which \(e\) is identity.
The assertions for \(d\) follow from (C.4a).
Chinese remaindering of \(N^{\rm ab}/6\) gives the
stated fibre-product description of \(G_6\).

The two nontriviality assertions can be decided without
selecting a compactness lift. We give both the precise
finite test and its independently run certificate.
The Bass--Serre graph of \(N\) has 576 edges, 288
vertices of type two and 192 of type three. Its cycle
rank is 97. Write \(M_p=H_1(N,\mathbb F_p)\).
For \(c\in\{r,s\}\), \(q\in C_Q(q(c))\), and any
word \(a\in\mathcal B\) over \(q\), the commutator
\([a,c]\) gives a cycle \(v_{a,c}\in M_p\).
Another lift \(na\), \(n\in M_p\), centralizes \(c\)
if and only if
\[
 v_{a,c}\in(c-1)M_p.
 \tag{C.7b}
\]
Changing the chosen word changes this cycle by an
 element of \((c-1)M_p\), so the test exhausts every
possible lift over \(q\).

The program \path{c_g3_frame_test.py} constructs
\(Q={\rm SL}_2(\mathbb Z/12)/\{\pm I\}\) from its
marked matrices, constructs this graph and an integral
spanning-tree cycle basis, and tests (C.7b) for every
quotient-centralizer element. Translation of edges
computes the exact peripheral action. Gaussian
elimination over \(\mathbb F_p\) gives
\[
\begin{array}{c|r|r|r|r}
 c&p&|C_Q(q(c))|&\operatorname{rank}(c-1)|_{M_p}
      &|\operatorname{im}(C_{G_p}(c)\to Q)|\\ \hline
 r&3&18&63&3\\
 r&2&18&68&18\\
 s&2&16&45&2\\
 s&3&16&52&16
\end{array}
\tag{C.7c}
\]
In the first row the image is exactly \(\langle q(r)\rangle\),
and excludes \([5I]\). Thus if \(e\) acted identically
on \(G_3\), its retained cubic frame over \([5I]\)
would give a centralizing lift excluded by (C.7b).
In the third row the image is exactly \(\langle q(s)\rangle\).
The two possible quartic-frame classes with positive
mod-four representative \(I\) and mod-three coordinate
\([i]\) have representatives
\[
 \begin{pmatrix}3&4\\8&3\end{pmatrix},\qquad
 \begin{pmatrix}3&8\\4&3\end{pmatrix}\pmod {12}.
\]
Neither belongs to that centralizer image. C.18 proves
that a normalized norm-one quartic frame has mod-four
representative \(\pm I\), so these are all possibilities.
If \(d\) acted identically on \(G_2\), its quartic
frame would contradict the third-row test.
The code saves every successful matrix and every
nonzero reduced defect in its JSON certificate;
its output verifies both exclusions. All matrices
are divided only by one global sign.
\end{proof}

The remaining finite action of a selected E.58 lift on
\(G_3\) has not been extracted from the compactness
construction; (C.7a) does not supply all its entries.
By C.18 there is one actual Galois element whose
\emph{entire adelic CM pair} equals that of \(de^{-1}\).
This improves the finite \((1,1)\)-bit comparison, but
does not identify its actions on \(G_6\) or the full
source. The implication
\(\Delta(de^{-1})=0\Rightarrow de^{-1}\in G_{\mathbb Q}\)
would require an image theorem for the full source,
which is substantially stronger than C.18's two CM
representations.

\begin{statement}{C.8}{The scope and exact-sequence law of the residual dimensions}{Proved}
The six numbers in R17 count specified small-character
components in their source statements. They are not the
whole dimensions of \(\mathcal J_n\), nor dimensions of
\((\mathcal J_n)_G\).
If \(\mathcal J_{n+1}=\ker(q_n|_{\mathcal J_n})\),
the exact law on a retained exact component is
\[
 \dim\mathcal J_{n+1}^{\rm sel}
  =\dim\mathcal J_n^{\rm sel}
    -\operatorname{rank}(q_n|_{\mathcal J_n^{\rm sel}}).
 \tag{C.8a}
\]
This identity supplies no recurrence for the ranks.
\end{statement}
\begin{proof}
Equation (C.8a) is rank-nullity, provided the map
preserves the specified component, or after applying
its exact idempotent projector.
Coinvariants in the modular characteristic are not
exact and cannot replace that component.
The source labels give
\[
\begin{array}{c|r|r}
 \hbox{LM7 row}&\hbox{selected dimension}
                    &\hbox{drop}\\ \hline
 D.60&353792&\\
 D.61&331776&22016\\
 D.62&325120&6656\\
 D.63&284160&40960\\
 D.64&282230&1930\\
 D.65&276478&5752
\end{array}
\tag{C.8b}
\]
We checked these labels and subtractions, not all
the large rank certificates. LM7 D.65 gives separately
\[
 \begin{aligned}
 \dim(\mathcal J_{64}/\mathcal J_{65})&=5752,\\
 \dim(\mathcal J_{64}/\mathcal J_{65})_G&=12,\\
 \dim(\mathcal J_{63}/\mathcal J_{65})_G&=26,
 \end{aligned}
 \tag{C.8c}
\]
and explicitly disclaims a formula for
\((\mathcal J_{65})_G\). It also retains a surjection
to the 1918-dimensional \(\ker\nu_4\).
The full D.63 dimension contains additional terms.
Thus replacing these selected dimensions by full
residual or obstruction dimensions changes the
source statement. This proves the scope correction
without independently certifying the numerical ranks.
\end{proof}

\begin{statement}{C.9}{Vanishing modulo the specified lower changes}{Open}
Task C.3 remains unresolved for the precise changes
allowed in LM7 D.61 and D.63--D.65.
A full lift forces zero obstruction after all changes
induced by its restrictions to stable quotients.
It has not been proved that those changes equal the
particular derivations, norm words, or homology
changes admitted by the recorded tower.
The requested dimension law and the arithmetic
boundary after the Lead 6 correction are also
undetermined.
\end{statement}

The existence implication is exact: if \(X\) is a
stable quotient for a full lift \(\Phi\), its induced
action is a witness to zero obstruction for its
lower restriction \(a_\Phi\).
Any obstruction quotient identifying \(a_\Phi\)
with the recorded lower representative \(a_0\)
therefore has zero class. The missing assertion is
that \(a_\Phi a_0^{-1}\) belongs to the stipulated
change module. We read the D.65 divided norms and
its restriction on fixed-homology corrections.
An extrapolation of (C.8b), or construction of a
further layer, would not prove this membership.

\begin{statement}{C.10}{The correct relative source for the excluded record}{Proved}
The marked quotient \(T_2=\mathcal B/D_8(K)\) in
LM7 D.32 is not a quotient of
\(\mathcal B/R_{\rm nil}(F)\).
Here \(K=\ker(\mathcal B\to Q_0)\), \(|Q_0|=18\),
is its rank-four source.
It is a quotient of \(\mathcal B/R_{\rm nil}(K)\)
and of \(\mathcal B/R_{\rm nil}(N)\).
The sphere quotient \(\mathcal B/D_9(F)\) in
LM7 D.45 is a quotient of
\(\mathcal B/R_{\rm nil}(F)\).
\end{statement}
\begin{proof}
The D.32 record retains \(Q\), with marked pure image
\[
 q(F)\simeq V_4\times{\rm SL}_2(\mathbb F_3).
 \tag{C.10a}
\]
This is not nilpotent: the order-three subgroup
generated by
\(\left(\begin{smallmatrix}1&1\\0&1\end{smallmatrix}\right)\)
is a nonnormal Sylow subgroup of
\({\rm SL}_2(\mathbb F_3)\).
The program \texttt{c\_cm\_obstruction.py} enumerates
the 96-element marked pure image modulo 12, its
24-element image at three, and a conjugate outside
this Sylow subgroup.
All finite quotients of \(F/R_{\rm nil}(F)\)
are nilpotent, proving the impossibility.
In contrast \(K/D_8(K)\) is a finite two-group,
as is its subgroup image of \(N\); both residuals
are killed. The final assertion follows because
\(F/D_9(F)\) is a finite two-group.
These are the actual quotient maps required in
the contrapositive lifting argument.
\end{proof}

\begin{statement}{C.11}{The central refinement detecting a finite obstruction}{Proved}
Suppose an action on a finite marked base, together
with the prescribed lower images of its three
peripheral conjugators, has no compatible peripheral
epimorphism lift to a specified refinement \(E\).
Construct \(D^\dagger\) as in C.1 from every
subgroup image and its invariant irreducibles.
Then that action has no compatible peripheral
epimorphism lift to \(D^\dagger\).
When the auxiliary source has nilpotent image
of \(U\), the action cannot be induced by a
positive peripheral automorphism of the whole
\(\Gamma_{U,\rm nil}\).
\end{statement}
\begin{proof}
A lift to \(D^\dagger\) would give the forbidden
lift to \(E\) by C.1.
A whole relative automorphism supplies such a
lift whenever the auxiliary source kills
\(R_{\rm nil}(U)\), as in C.2.
This locates the detecting central layer by an
explicit finite construction; it does not require
guessing an individual obstructing irreducible.
\end{proof}

\begin{statement}{C.12}{Exact central detectors and the cusp normalization gap}{Open}
The normalized finite exclusions of LM7 D.32 and D.45
are independently certified. Each has an explicit
central \(C_2\)-refinement \(D^\dagger\to D\) detecting
its recorded action with cusp conjugator one.
Consequently neither recorded action is induced by
a \emph{cusp-normalized} peripheral automorphism of
its correctly typed relative source in C.10.

The extension of this conclusion to an arbitrary
cusp conjugator is not supplied by those exclusions.
In fact the sphere action of LM7 D.45 \emph{does}
lift to a peripheral automorphism of
\(\mathcal B/D_{10}(F)\) when its cusp may be conjugated,
while inducing exactly the recorded action on
\(\mathcal B/D_9(F)\).
This refutes the unrestricted finite nonlifting
inference. Whether the recorded action extends
with arbitrary cusp conjugator to the \emph{whole}
relative completion remains open here; thus task
C.4 has this explicit distinction of scope.
\end{statement}
\begin{proof}[Proof of the normalized assertions and finite counterexample]
Write \(D_n(U)\) for the two-Zassenhaus filtration.
For the D.32 case take \((U,n)=(K,8)\), and for the
sphere D.45 case take \((U,n)=(F,9)\). Set
\[
 E=\mathcal B/D_{n+1}(U),\quad D=\mathcal B/D_n(U),\quad
 L_n=D_n(U)/D_{n+1}(U).
 \tag{C.12a}
\]
The layers have dimensions \(8230\) and \(56\).
They are the actual homogeneous restricted Lie layers:
the powers of Lyndon commutators of degrees
\(\lvert w\rvert 2^j=n\) form a basis of the
Zassenhaus quotient \citep[Proposition 4.4 and its
proof; Corollary 3.4]{cEfrat2023}.
Thus the Magnus calculations below concern exact
finite groups, not only a chosen subspace of their
associated graded objects.

\emph{What was rerun.}
The wrapper \path{c_replay_obstructions.py} copies
the attached code into an isolated work directory,
regenerates the omitted JSON inputs, and reevaluates
the literal recorded words and all the relevant
centralizer and multiplier corrections. Its new
outputs and certificates have prefix \path{c_replay_}.
The recorded D.39 integer gauge \(-438\) is an input
defining the action; its earlier homology search
is not silently claimed as rerun. Source hashes
and this distinction are retained in the manifest.

For D.32 the \(4{,}903{,}248\)-letter cusp word has
zero Magnus coefficients through degree seven,
and its degree-eight target has support \(32296\).
The central correction map is
\[
 (b,a)\longmapsto(1-s)b+s(1-r)a.
 \tag{C.12b}
\]
Its rank is \(7753\). The cusp orders in \(D,E\)
are \(48,96\), respectively, so all compatible
multiplier changes add exactly one possible
column, the degree-eight term of \(\tau^{48}\).
The rank becomes \(7754\).
The complete lower \(s\)-centralizer in
\(N/D_8(K)\) has \(1666\) generators, with leading
multiplicities \(37,102,348,1179\) in degrees
four through seven. Adding every one of their
defects gives rank \(7777\), still excluding the target.

Here completeness is a group statement. In a
successive central extension with layer \(L_j\),
the commutation defect defines a homomorphism
from the lower centralizer to \(L_j/(1-s)L_j\).
Its kernel is the image of the upper centralizer,
and the new kernel is \(L_j^s\).
Descending elimination in a full polycyclic sequence
constructs a sequence for this kernel. The program
checks the actual word substitutions and independent
leading coordinates; its ordered products therefore
have at least the count dictated by the exact
kernel formula, and that formula gives the matching
upper bound. Starting from order \(2^6\), the
successive centralizer orders have base-two logarithms
\(16,50,152,494,1666\).
The final depth bound makes the last group elementary
abelian. This proves that its generators exhaust
all lower changes, not merely their linearizations.
The inclusion \(D_8(K)\subset[N,N]N^2\) is also
verified: the ordinary commutator images have ranks
\(87,67,43,19,4,0\), and the remaining restricted
powers vanish in \(N/[N,N]N^2\).
Thus \(D\) retains \(G_2\), and C.7's independent
centralizer calculation forces every normalized
lower \(s\)-frame difference into \(N/D_8(K)\).
The order-three centralizer lifts across a two-group
kernel by coprime complement conjugacy, and supplies
no omitted class.

For D.45 the replay evaluates the
\(12{,}921{,}196\)-letter cusp word at every starting
vertex of the two six-edge graphs, with \(108\)
independent tuple-monomial checks of the exact
arithmetic. The sphere target vanishes through
degree eight. The \(112\) columns of (C.12b) have
rank \(47\). The entire lower sphere
\(s\)-centralizer has order \(2^{38}\), constructed
through the ranks
\(1,2,3,5,8,13,22,38\); all its next-layer defects
are already in that rank-\(47\) space.
The same exact connecting-map induction proves
completeness, with actual free words for every
generator. The order-three correction is again
coprime. Both the lower and upper cusp orders
are \(32\), so every multiplier preserving the
lower action is invisible at this transition.
As a control, the elliptic rank-two case has
central rank \(38\); its \(50\) lower-centralizer
generators enlarge this to \(56\), and it does lift.
We do not extend the sphere exclusion to that case.

\emph{The actual central covers.}
Let \(\Lambda_{32}:L_8\to\mathbb F_2\) be the
722-coefficient functional written exactly in
\path{c_d32_central_character.hex}, using the
rank-four basis and coefficient order specified
in the regenerated certificate. For the sphere,
in its fixed free basis write \(0=X,1=Y\), and set
\[
 \begin{split}
 \Lambda_{45}(z)={}&[000011011]z+[000010111]z\\
                  &+[001001111]z+[000101111]z.
 \end{split}
 \tag{C.12c}
\]
These are little-endian tensor positions
\(432,464,484,488\). Its four target coefficients
are \((0,0,1,0)\).
Both functionals annihilate every central correction,
every complete lower \(s\)-centralizer defect,
and every permitted multiplier column, while taking
value one on the target. The independently run
\path{c_central_refinements.py} reconstructs the
restricted bases and the rank-four quotient graph,
and checks
\[
 \Lambda(sz)=\Lambda(rz)=\Lambda(z)\qquad(z\in L_n)
 \tag{C.12d}
\]
on all \(8230\), respectively \(56\), basis elements.
This invariance also follows abstractly from
annihilation of both summands of (C.12b): first
\(\Lambda s=\Lambda\), then \(\Lambda r=\Lambda\).

Define the closed normal subgroup and its actual
source quotient by
\[
 J_\Lambda=\{u\in D_n(U):
               \Lambda(\operatorname{Magnus}_n(u))=0\},\qquad
 D^\dagger=\mathcal B/J_\Lambda.
 \tag{C.12e}
\]
Equation (C.12d) proves normality in all of
\(\mathcal B\), and the nonzero target proves
\[
 1\longrightarrow C_2\longrightarrow D^\dagger
       \longrightarrow D\longrightarrow1
 \tag{C.12f}
\]
is a central extension. Its order is
\(18\cdot2^{3339}\) in the rank-four case and
\(6\cdot2^{85}\) in the sphere case.
This is precisely the one-dimensional invariant
character \(\rho(z)=(-1)^{\Lambda(z)}\) in C.1's
count. Indeed the image of \(E\) in its scalar
cover, by \(e\mapsto[(e,1)]\), is
\(E/\ker\rho=D^\dagger\); the original source
map gives the required marked epimorphism.
Thus an explicit character, and not only the
existence of an unspecified central cover, locates
the obstruction asked for in task C.4.

With cusp conjugator one, a hypothetical lift
would change the target only by the columns just
exhausted. Applying \(\Lambda\) would give
\(1=0\). This proves normalized nonlifting even
to \(D^\dagger\). Since its image of \(U\) is a
finite two-group, the quotient map factors through
the corresponding whole relative nilpotent source.
A cusp-normalized peripheral automorphism of that
source inducing the recorded lower action would
therefore give the forbidden lift.

\emph{Why the cusp hypothesis cannot be discarded.}
If a peripheral automorphism merely induces the
normalized action on \(D\), its actual cusp frame
may project to a nontrivial element of \(C_D(\tau)\).
Those possibilities are not among the preceding
\(s\)-centralizer changes.
The separate program \path{c_cusp_centralizer_scope.py}
constructs the complete subgroup
\(C_{F/D_9(F)}(\tau)\), using the actual nonlinear
substitution by \(\tau=sr\). The same exact
central-extension argument gives logarithmic orders
\[
 1,3,3,6,6,9,12,23
 \quad\text{in }F/D_j(F),\quad j=2,\ldots,9.
 \tag{C.12g}
\]
The final leading degrees are \(1,2,4,8\), with
multiplicities \(1,1,1,20\).
The \(23\) next-layer cusp defects enlarge the
rank-\(47\) correction space to the entire
56-dimensional layer. In particular three of
them have \(\Lambda_{45}=1\).
A cusp frame outside the prescribed fibre can
therefore cancel the scalar obstruction.

More concretely, the saved solution gives actual
free words \(b,a\in D_9(F)\) of lengths \(3626,1292\)
and \(v\in F\) of length \(834\). Put
\(\widetilde\beta=b\beta\),
\(\widetilde A=aA\), using the exact recorded
conjugators, and define
\[
 \widetilde\alpha(s)=\widetilde\beta s\widetilde\beta^{-1},
 \qquad
 \widetilde\alpha(r)=\widetilde A r\widetilde A^{-1}.
 \tag{C.12h}
\]
Full noncommutative multiplication, rather than a
comparison only of leading terms, verifies in
\(\mathcal B/D_{10}(F)\) that
\[
 \widetilde\alpha(s)\widetilde\alpha(r)\tau^{-1}
        =v\tau v^{-1}\tau^{-1}.
 \tag{C.12i}
\]
All correction words and both equal polynomials
are saved in the JSON certificate. The words
\(b,a\) are trivial in the lower base, so this
is exactly the original lower action.
The torsion relations hold because the two
images are actual conjugates. The induced map
on \(F/D_2(F)\) is the recorded identity, and
the quotient \(S_3\) action is identity; the
finite Frattini theorem makes the endomorphism
surjective, hence an automorphism.
This is the asserted finite counterexample to
removing cusp normalization. It does not assert
an extension through every subsequent layer.
\end{proof}

\subsection{The pentagon and the arithmetic commutant}

\begin{statement}{C.13}{The extra weight-nine direction fails the pentagon}{Proved}
Assume the existence assertion R16, specifically
LM7\_PUSH D.2 and D.52: a full multiplier-one
gentle pair has first pro-\(\ell\) term in weight
nine, for some \(\ell\ge5\), with metabelian
pair coordinate
\[
 P_9(X,Y)[x,y],\qquad P_9=X^2Y^2(X+Y)^3.
 \tag{C.13a}
\]
Every full lift retaining this leading term
fails the profinite pentagon.
Under exactly this existence hypothesis,
\(\widehat{\rm GT}\subsetneq\widehat{\rm GT}_0\).
\end{statement}
\begin{proof}
Map the free profinite group to its pro-\(\ell\)
quotient, then to its rational nilpotent quotients,
and take logarithms and the first nonzero
homogeneous term. For a full-GT element this
term lies in \(\mathfrak{grt}_1(\mathbb Q_\ell)\).
We need only this direction of comparison.
Expand the two-cycle, three-cycle and five-cycle
equations in that first weight. Products become
sums. The initial pure-braid relations are
\([t_{ij},t_{kl}]=0\) for four distinct indices
and \([t_{ij},t_{ik}+t_{jk}]=0\).
Successively forgetting a strand has a free
kernel and trivial action on its abelianization;
induction gives the associated graded Lie
algebra with these relations. The expanded
pentagon is therefore the infinitesimal
pentagon in that algebra.

Furusho's inclusion
\(\mathfrak{grt}_1\subset\mathfrak{dmr}_0\)
and gamma factorization apply
\citep[Corollary 0.3, Theorem 0.5, \S5]{cFurusho2011}.
Here is the coordinate comparison.
Write \(v=v_xx+v_yy\) in the completed
associative algebra and set \(M(v)=(v_yy)^{\rm ab}\).
Direct multiplication gives
\[
 \begin{aligned}
 M([x,y])&=XY,\\
 M([x,v])&=XM(v),\qquad M([y,v])=YM(v)\quad(v\in L').
 \end{aligned}
 \tag{C.13b}
\]
Thus \(M\) kills \(L''\) and sends
\(p(X,Y)[x,y]\) to \(XYp(X,Y)\).
Differentiating the gamma factorization gives
\(M(v)=\gamma(X)+\gamma(Y)-\gamma(X+Y)\).
In degree nine this is a scalar multiple of
\(X^9+Y^9-(X+Y)^9\). Multiplication by \(XY\)
is injective, so the full-GT metabelian
weight-nine image lies in
\[
 \mathbb Q_\ell G_9,\qquad
 G_9=\frac{(X+Y)^9-X^9-Y^9}{XY}.
 \tag{C.13c}
\]
The coefficient of \(X^7\) in \(P_9\) is zero
and in \(G_9\) is 9; the coefficient of
\(X^5Y^2\) in \(P_9\) is one.
They are independent over \(\mathbb Q_\ell\).
The exact coefficient minor is 9, checked by
\texttt{c\_cm\_obstruction.py}.
Hence (C.13a) violates a necessary consequence
of the pentagon, for every lift retaining it.
\end{proof}

The dependence is exact. Existence is the
unrefereed portion of LM7\_PUSH D.2/D.52,
with the LM7\_PUSH D.20, D.21, D.51 and D.52
inputs listed in R16. The D.50 lifting step
has an attached favourable referee verdict
and is reconstructed in C.1--C.2.
The gamma upper bound is LM7 D.21's comparison,
with its primary theorem opened and its
coordinate map supplied above.
R16's further input audit belongs to the
parallel prompt specified by the user.
We do not state an unconditional new
strict-containment theorem.

\begin{statement}{C.14}{R16 lies outside the arithmetic commutant}{Proved}
A multiplier-one metabelian peripheral outer
action with a first nonzero rational
pro-\(\ell\) derivation of positive weight
cannot commute with \(G_{\mathbb Q}\).
Thus the R16 weight-five and weight-nine
classes are not the arithmetic-commuting
classes of R18.
Moreover LM6 A.49, A.56 and A.58 exclude
Galois-commuting third-derived lifts,
not arbitrary peripheral third-derived lifts.
\end{statement}
\begin{proof}
On the graded pro-\(\ell\) tripod Lie algebra,
\(g\) acts in degree \(j\) by
\(\chi_\ell(g)^j\), because both degree-one
generators have cyclotomic action.
Conjugation therefore multiplies a derivation
shifting degree by \(n\) by \(\chi_\ell(g)^n\).
The assertion survives quotienting by inner
derivations and the second derived ideal.
Commutativity of a nonzero leading class \(v\)
with every \(g\) would imply
\((\chi_\ell(g)^n-1)v=0\) for all \(g\).
Cyclotomic surjectivity supplies a value
of infinite order, a contradiction over
\(\mathbb Q_\ell\).

The exact source statement LM6 A.58 \citep{cLM6} is
\[
 \operatorname{im}\left(
 C_{\operatorname{Out}_{\rm per}(F/F''')}(G_K)
       \longrightarrow\operatorname{Out}(F/F'')\right)=S_3.
 \tag{C.14a}
\]
Its source is a centralizer. LM6 A.49 says
explicitly that its unit family has no
commuting third-derived lift.
That family is trivial on every maximal
pro-prime quotient, whereas R16 has the
nonzero coordinate just considered.
The two claims concern different subsets
and different lifting conditions.
Removing the centralizer condition from
R18 strengthens the source without proof.
We verify this scope and the compatibility,
not the entire proof of (C.14a).
\end{proof}

\begin{statement}{C.15}{The CM-parity restriction in task C.6}{Open}
It is not proved here whether
\(\Delta|_{\widehat{\rm GT}}=0\),
or whether the full beta and gamma
compatibilities evaluated on CM coordinates
force it. Rational pro-prime gamma
constraints alone do not establish it.
\end{statement}

The E.58 element reconstructed in C.18
witnesses the last assertion: its pure
pronilpotent action is identity, so all
rational pro-prime gamma coordinates vanish,
while \(\Delta=1\).
Full profinite beta data also contain
prime-to-coefficient deck coordinates, lost
under that comparison.
We opened \citet[Theorem 8.1 and Remark 8.6]{cEnriquez2007},
which proves \(\widehat{\rm GT}_K=\widehat{\rm GT}\)
and cyclic-cover compatibility.
It does not itself identify this CM parity.
The remaining step is an explicit comparison
from the full distribution identities to
the integral CM restrictions and their
common quadratic extension. A graded
necessary condition cannot reverse that map.

The precise literature distinction is visible in
\citet{cNakamuraTsunogaiYasuda2010}.
Their Proposition 4.4 computes the Mordell triangle-cover
Tate action by a beta specialization for full GT.
Their Proposition 4.13 and Theorem 4.16, however, identify
the braid path and the equianharmonic frame only for
\(\sigma\in G_{\mathbb Q}\); the latter proof uses
Proposition 4.11's arithmetic path formula.
Applying that frame identity to every full-GT element
would add the very comparison still needed here.

C.26 proves that the folded two-dimensional four-strand
pentagon test is identically satisfied by all gentle pairs.
Its failure to detect parity is exact, at every modulus.
The distinct-endpoint tests recorded in C.27 have their
finite necessary-condition scope stated separately.
C.28 gives a different structural exclusion: a
perfect common braid image is independent of every
finite solvable source. Its explicit three-dimensional
modulo-eleven example therefore cannot remove any
prescribed \(G_6\) conjugator fibre. Neither result
extends to the full profinite pentagon.

\subsection{Lead 7: the exact integral CM image}

\begin{statement}{C.16}{The whole norm-one product in Lead 7 is too large}{Refuted}
Read without the necessary modulo-two conditions,
Lead 7 is false. For the normalized projective adelic CM
coordinates of LM7\_PUSH E.65, the image
of \(\mathcal H\) is not the whole product
of the two integral projective norm-one groups.
Its Gaussian coordinate is identity modulo
2, whereas the norm-one element \([i]\)
is not. Thus the adelic pair \((1,[i])\)
in that whole product has no preimage.
\end{statement}
\begin{proof}
Set
\[
 U=\begin{pmatrix}1&1\\0&1\end{pmatrix},\quad
 L=\begin{pmatrix}1&0\\-1&1\end{pmatrix},\quad
 S=ULU=\begin{pmatrix}0&1\\-1&0\end{pmatrix}.
 \tag{C.16a}
\]
For \(a\in\mathcal H\), its normalized
pure \(s\)-frame \(\beta_a\) belongs to \(F\).
With normalized upper-unipotent congruence
intertwiner \(V_a\), the Gaussian coordinate
is \(W_a=q(\beta_a)^{-1}V_a\), projectivized.
The first factor is identity modulo 2.
The normalized action induces identity on
the \(S_3\) quotient, so \(V_a\bmod2\)
commutes with \(U\bmod2,L\bmod2\).
The simultaneous centralizer among upper
shears is the identity, by direct
multiplication. The program
\path{c_cm_obstruction.py} also
enumerates it exactly. Hence \(W_a\bmod2=I\).

The Gaussian norm-one group is
\[
 T_i^1(\widehat{\mathbb Z})
 =\{aI+bS:\ a,b\in\widehat{\mathbb Z},\
                         a^2+b^2=1\}.
 \tag{C.16b}
\]
It contains \(S\), the element \(i\).
Its reduction
\(\left(\begin{smallmatrix}0&1\\1&0\end{smallmatrix}\right)\)
is nonscalar over \(\mathbb F_2\); its
projective class cannot equal \(W_a\).
This finite obstruction does not conflict
with C.4's three-primary coordinate \([i]\),
which is not a claim that the entire
adelic coordinate is the global unit \(i\).
\end{proof}

The same argument applies to the cubic frame, which
also belongs to \(F\). Both CM coordinates are therefore
identity modulo two. The quotient by signs in this
paragraph is always by one global \(\{\pm I\}\).
C.18 specifies these necessary conditions and proves
surjectivity onto the resulting product. Thus, if the
``norm-one groups'' in Lead 7 were intended to include
these conditions, C.18 proves that precise interpretation.

\begin{statement}{C.17}{The field-intersection sublemma in Lead 7}{Proved}
In a fixed algebraic closure,
\[
 \mathbb Q(i)^{\rm ab}
   \cap\mathbb Q(\sqrt{-3})^{\rm ab}
       =\mathbb Q^{\rm ab}(3^{1/4}).
 \tag{C.17a}
\]
This does not assert a full integral
image for either CM representation.
\end{statement}
\begin{proof}
Call the intersection \(L\).
It is Galois over \(\mathbb Q\) and
contains \(\mathbb Q^{\rm ab}\).
Then \(G=\operatorname{Gal}(L/\mathbb Q)\)
has distinct abelian index-two subgroups
\(H_1,H_2\), fixing the two quadratic fields.
Their intersection \(H_0\) is central
since they generate \(G\).
Choose \(a\in H_1\setminus H_0\),
\(b\in H_2\setminus H_0\).
They, with \(H_0\), generate \(G\), and
\(a^2,b^2\in H_0\).
The closed commutator subgroup is
generated by the central \([a,b]\),
whose square is \([a^2,b]=1\).
Its fixed field is precisely
\(\mathbb Q^{\rm ab}\), so
\([L:\mathbb Q^{\rm ab}]\le2\).

Let \(\alpha=3^{1/4}\).
The splitting field \(\mathbb Q(i,\alpha)\)
has dihedral group \(D_8\), with
\(r(\alpha)=i\alpha,r(i)=i\) and
\(s(\alpha)=\alpha,s(i)=-i\).
The subgroup fixing \(\mathbb Q(i)\)
is \(\langle r\rangle\simeq C_4\).
Since \(\sqrt{-3}=i\alpha^2\), that
fixing \(\mathbb Q(\sqrt{-3})\) is
\(\langle r^2,rs\rangle\simeq C_2^2\).
Both are abelian, so this field is in
\(L\). The program
\path{c_cm_obstruction.py}
checks these exact subgroup actions.
But \(\mathbb Q(\alpha)/\mathbb Q\)
is not Galois, and thus is not a
subfield of an abelian extension.
The degree bound is attained.
\end{proof}

\begin{statement}{C.18}{Surjectivity onto the corrected CM product and its sole parity}{Proved}
Put \(K_3=\mathbb Q(\sqrt{-3})\), \(K_4=\mathbb Q(i)\),
\(\mathcal O_3=\mathbb Z[\zeta_3]\),
\(\mathcal O_4=\mathbb Z[i]\), and define profinite groups
\[
 \begin{split}
 \mathcal U_j&=
  \{u\in\widehat{\mathcal O}_j^{\times}:u\equiv1\pmod2\}
                                    /\{1,-1\},\\
 \mathcal T_j&=\ker\bigl(N:\mathcal U_j
                                  \to\widehat{\mathbb Z}^{\times}\bigr),
                  \qquad j=3,4 .
 \end{split}
 \tag{C.18a}
\]
Each denominator is the single global group of order
two. The continuous CM homomorphism is onto:
\[
 \Theta:\mathcal H\longrightarrow\mathcal T_3\times\mathcal T_4.
 \tag{C.18b}
\]
Each class in \(\mathcal T_j\) has a unique
representative congruent to \(I\pmod4\).
For these representatives set
\[
 \begin{split}
 \varepsilon_3(T)&=0\text{ or }1
    \quad\text{according as }T(1,-1)=+(1,-1)
                       \text{ or }-(1,-1)\pmod3,\\
 \varepsilon_4(W)&=a\pmod2\qquad(W\bmod3=S^a).
 \end{split}
 \tag{C.18c}
\]
The exact joint image of \(G_{\mathbb Q^{\rm ab}}\) is
\[
 \mathcal Z=\ker(\varepsilon_3+\varepsilon_4)
                \subset\mathcal T_3\times\mathcal T_4,
 \qquad [\mathcal T_3\times\mathcal T_4:\mathcal Z]=2.
 \tag{C.18d}
\]
The parity \(\delta=(\varepsilon_3+\varepsilon_4)\Theta\)
has a unique continuous extension \(\Delta\) to the
full gentle group that vanishes on \(G_{\mathbb Q}\).
It is the character of LM7\_PUSH E.65 (R14).
Moreover \(\Theta(\ker\Delta\cap\mathcal H)=\mathcal Z\).
For the full pair \(\widetilde\Theta\) with values in
the two projective CM normalizers, we also have
\[
 \widetilde\Theta(\ker\Delta)
           =\widetilde\Theta(G_{\mathbb Q}).
 \tag{C.18n}
\]
Thus \(\Delta\) is exactly the obstruction in these
two CM coordinates. It is not asserted to be the only
obstruction to an element of \(\widehat{\rm GT}_0\)
being Galois.
\end{statement}
\begin{proof}
We first check the target and the homomorphism, then
identify the actual arithmetic images, and finally
construct a gentle element in the other coset.

The signed braid representation sends the two braid
generators to \(U,L\) of (C.16a). For a gentle element
\(a\) its precomposition sends the first to
\(U^{\lambda(a)}\) and the second to a trace-two
matrix \(Y\) conjugate to \(L^{\lambda(a)}\), satisfying
the braid relation. This uses precisely the two
hexagons and the actual automorphism, as in the modular
description of \citet{cHarbaterSchneps2000}; no pentagon
is used. Over the domain \(\mathbb Z_p\), the lower
left entry \(c\) of that braid relation satisfies
\(c(1+\lambda c)=0\).
If \(c=0\), trace and determinant force
\(Y=U^\lambda\), contradicting the nonabelian image
of this precomposition. Otherwise
\[
 Y=\begin{pmatrix}1-u&\lambda u^2\\
                   -1/\lambda&1+u\end{pmatrix},\qquad
 V_a=\begin{pmatrix}\lambda&\lambda u\\0&1\end{pmatrix}
 \tag{C.18e}
\]
gives its covariance. These calculations also prove
uniqueness and are the rigidity calculation of
\citet[Lemmas 4.1--4.3]{cMarin2008}.
Taking the product over all primes gives the continuous
normalized adelic matrix, with
\(q\Phi_a=\operatorname{Ad}(V_a)q\) and
\(V_{ab}=V_aV_b\).

For \(a\in\mathcal H\), let \(A_a,\beta_a\in F\)
be its unique pure frames. They depend continuously on
\(a\): conjugation from the compact pure group to the
given conjugacy class is a continuous bijection onto
a Hausdorff space. Their uniqueness gives
\[
 A_{ab}=\Phi_a(A_b)A_a,\qquad
 \beta_{ab}=\Phi_a(\beta_b)\beta_a.
 \tag{C.18f}
\]
Consequently
\(T_a=q(A_a)^{-1}V_a\) and
\(W_a=q(\beta_a)^{-1}V_a\), modulo the global sign,
are continuous homomorphisms.
Both are identity modulo two by C.16. They commute
projectively with \(R_+=UL\) and \(S\), respectively.
The first commutation is genuine, since
\(\operatorname{tr}R_+=1\) excludes changing its sign.
The second is genuine too: an anticommuting matrix
with \(S\) has shape
\(\left(\begin{smallmatrix}a&b\\b&-a\end{smallmatrix}\right)\).
If it is identity modulo two its determinant is
\(-a^2-b^2\equiv-1\pmod4\), contradicting determinant
one. The two centralizer rings are exactly
\(\widehat{\mathbb Z}[R_+]\simeq\widehat{\mathcal O}_3\)
and \(\widehat{\mathbb Z}[S]\simeq\widehat{\mathcal O}_4\);
this follows by solving their two-by-two commuting
equations, or by the cyclic vectors of these matrices.
Determinant is the quadratic norm. Thus \(\Theta\)
has precisely the target asserted in (C.18b).

We next prove the individual arithmetic images, including
their exact integral subgroup. Both CM rings are
Euclidean, hence have class number one. CM reciprocity
identifies the Weber-function torsion fields with the
ray class fields, and their Galois actions with the
quotients of the finite CM unit groups
\citep[Corollaries 3.4 and 3.7, Theorem 4.1]{cLozanoRobledo2022}.
The compatible adelic form is also the torsion diagram
of \citet[Theorem 3.10]{cMilne2007}.
For a direct image argument, take the Weber functions
\(x^3\) on \(E_3\) and \(x^2\) on \(E_4\).
CM endomorphisms are defined over \(K_j\), and each
Tate module is free of rank one over its completed CM
order. Thus the Tate image is contained in the
commutative group \(\widehat{\mathcal O}_j^{\times}\).
At every finite level, the kernel of its action on all
Weber values is its intersection with the image of
\(\mathcal O_j^{\times}\). To check this last assertion,
evaluate on a generator of the rank-one
\(\mathcal O_j/n\)-module: equal Weber values differ
by one CM automorphism, which then acts identically
on the whole module. Reciprocity gives the full
quotient \((\mathcal O_j/n)^\times/\operatorname{im}
(\mathcal O_j^\times)\). Compactness over all levels
therefore gives the full adelic quotient
\(\widehat{\mathcal O}_j^\times/\mathcal O_j^\times\),
as an actual image, not only an abstract group isomorphism.

Reduction modulo two maps
\(\mathcal O_j^\times/\{\pm1\}\) isomorphically to
\((\mathcal O_j/2)^\times\): these groups have order
three for \(j=3\), and order two for \(j=4\).
Every unit can therefore be uniquely multiplied by
such a root of unity to become one modulo two, modulo
the global sign. Hence the natural map
\[
 \mathcal U_j\xrightarrow{\ \sim\ }
       \widehat{\mathcal O}_j^\times/\mathcal O_j^\times
 \tag{C.18g}
\]
is an isomorphism. Since \(E_j[2]\) is defined over
\(K_j\), its projective Tate image is contained in
\(\mathcal U_j\). Surjectivity through the isomorphism
(C.18g) proves that its image is exactly \(\mathcal U_j\).
Moreover its field is exactly \(K_j^{\rm ab}\): it
contains all the Weber ray class fields, and is
contained in \(K_j^{\rm ab}\) by CM linearity.
Its determinant is the cyclotomic character by the
Weil pairing. Restriction to \(G_{\mathbb Q^{\rm ab}}\)
therefore has image exactly \(\mathcal T_j\).
The based comparison C.5 identifies these with the
two arithmetic CM coordinates under consideration.

C.17 now determines their joint image by ordinary
Galois theory. It is the fibre product of the two
individual groups over
\[
 \operatorname{Gal}\bigl(
  \mathbb Q^{\rm ab}(3^{1/4})/\mathbb Q^{\rm ab}\bigr)=C_2.
 \tag{C.18h}
\]
We identify its restrictions with (C.18c), to retain
the sign convention. A norm-one Eisenstein unit
\(a+b\zeta_3\equiv1\pmod2\) has \(b\equiv0\pmod4\)
by \(a^2-ab+b^2\equiv1\pmod4\), hence is
\(\pm1\pmod4\). A Gaussian unit with \(a\) odd and
\(b\) even has \(b\equiv0\pmod4\), since otherwise
\(a^2+b^2\equiv5\pmod8\). It too is \(\pm1\pmod4\).
Exactly one global sign makes either representative
identity modulo four. Reduction of these representatives
is a homomorphism at three, where their centralizers
are the cyclic groups
\[
 C_{{\rm SL}_2(\mathbb F_3)}(-R_+)
       =\langle-I\rangle\times\langle-R_+\rangle=C_6,
 \qquad C_{{\rm SL}_2(\mathbb F_3)}(S)=\langle S\rangle=C_4.
 \tag{C.18i}
\]
These equations follow by direct matrix multiplication.
The rational three-torsion line on \(E_3\) is
\(\ker(-R_+-I)=\mathbb F_3(1,-1)\).
Changing the positive-modulo-four representative by
its sign on this line gives the actual arithmetic
sign fixing \((0,1)\); its action on \(E_3[4]\) is
then \((-I)^{\varepsilon_3}\). C.3's division-point
formulas identify this with the action
\(r_0\mapsto(-1)^{\varepsilon_3}r_0\).
On the Gaussian side \(S^a\) multiplies the
three-torsion \(x\)-coordinates by \((-1)^a\);
the formulas in C.3 identify that with
\(r_0\mapsto(-1)^{\varepsilon_4}r_0\).
Since \(r_0=3^{1/4}/\sqrt2\) and \(\sqrt2\in\mathbb Q^{\rm ab}\),
these are precisely the two restrictions in (C.18h).
This proves (C.18d) and identifies its parity with
LM7\_PUSH E.65's \(\delta\).

We check the extension, including invariance, to
reconstruct R14. For a general gentle element \(a\),
put \(\epsilon_a=\lambda(a)\bmod3\in\{1,-1\}\) and
\(b_a=(1-\epsilon_a)/2\). Choose the unique frames
\[
 A_a\in Fs^{b_a},\quad\beta_a\in F,\qquad
 \Phi_a(r)=A_a r^{\epsilon_a}A_a^{-1},\quad
 \Phi_a(s)=\beta_a s\beta_a^{-1}.
 \tag{C.18l}
\]
Existence follows by right multiplication of any frame
by its cyclic inertia generator: the action on
\(S_3\) is identity, so these are exactly the required
cosets. The torsion centralizer theorem proves uniqueness.
The product formulas (C.18f) hold with these cosets
because \(b_{ab}=b_a+b_b\pmod2\) and the action on
\(S_3\) is identity. Thus the projective matrices
\(\mathcal T_a=q(A_a)^{-1}V_a\) and
\(\mathcal W_a=q(\beta_a)^{-1}V_a\) are homomorphisms
on the full gentle group. They normalize the cyclic
groups generated by \(-R_+\) and \(S\), respectively.
The order-three statement is genuine by its trace;
the order-four statement can send \(S\) to \(-S\),
which still normalizes \(\langle S\rangle\).

Consequently conjugation by \(a\) acts on each
multiplier-one coordinate by these matrices.
Its positive-modulo-four representative remains
positive under conjugation. At three the cubic
normalizer preserves the fixed line of \(-R_+\),
so it preserves its scalar sign; the Gaussian
normalizer acts on \(C_4\) by an automorphism,
so preserves exponent parity. Hence \(\delta\)
is invariant under the full gentle group. It vanishes
on \(G_{\mathbb Q^{\rm ab}}\) by (C.18d).
For \(a\) choose \(\sigma\in G_{\mathbb Q}\) with
\(\chi(\sigma)=\lambda(a)\), and set
\[
 \Delta(a)=\delta(a\sigma^{-1}).
 \tag{C.18m}
\]
The cyclotomic character is onto, so such a choice
exists. Two choices differ in \(G_{\mathbb Q^{\rm ab}}\),
which proves independence. Invariance proves the
homomorphism law. Finally the continuous surjection
\(\mathcal H\times G_{\mathbb Q}\to\widehat{\rm GT}_0\)
is a quotient map, being a map from a compact space
onto a Hausdorff space. The continuous function
\((a,\sigma)\mapsto\delta(a)\) factors through it,
proving continuity without requiring a continuous
choice of \(\sigma\). The same product decomposition
proves uniqueness. This completes the independent
character reconstruction.

Finally we reconstruct the required nonzero coset of
LM7\_PUSH E.58, rather than using that row as authority.
Put \(H_0=F'F^3\), so \([\mathcal B:H_0]=54\).
In the actual signed pure congruence image modulo 12,
\[
 q(F)=V_4\times{
m SL}_2(\mathbb F_3),\qquad
 q(H_0)=V_4\times Q_8,
 \qquad q([x,y]^2)=5I.
 \tag{C.18j}
\]
The exact enumeration in
\path{c_corrected_cm_image.py} checks the orders
\(96,32\), the commutator square, and both finite
parity maps. In particular \(z_0=[5I]\) is a central
element of order two in \(q(H_0)\).
Let \(J\) be the actual joint image of \(\mathcal B\)
in \(Q\times\mathcal B/H_0\). It is a quotient of
\(\mathcal B/O^2(H_0)\), since its kernel over
\(\mathcal B/H_0\) is the two-group \(q(H_0)\).
The element \((z_0,1)\) belongs to \(J\), being the
image of \([x,y]^2\), and centralizes its order-three
generator. Coprime complement conjugacy, as proved in
C.4, lifts it to
\[
 c_2\in C_{\mathcal B/O^2(H_0)}(r)
                   \cap H_0^{(2)}.
\]
The tuple \(c=(c_2,1,1,\ldots)\) is consequently an
actual centralizer element in
\(\mathcal B/R_{\rm nil}(H_0)\), with common base
coordinate one. Give the identity automorphism of
that whole quotient the three frames \((1,c,1)\)
at \((s,r,\tau^{-1})\). C.2 supplies a full peripheral
automorphism retaining these exact frames and fixing
\(\tau\) literally. It induces identity on the pure
pronilpotent quotient, hence its normalized \(f\)
belongs to \(F'\); it is a multiplier-one gentle
element \(e\), by the two torsion relations as in C.4.
Its congruence action is identity and its frame
coordinates are \(q(A_e)=z_0,q(\beta_e)=1\).
The normalized shear commuting with both \(U,L\)
is identity, so
\[
 T_e\equiv5I\pmod{12},\qquad W_e\equiv I\pmod{12},
 \qquad (\varepsilon_3,\varepsilon_4)(\Theta(e))=(1,0).
 \tag{C.18k}
\]
The image of \(\mathcal H\) contains the index-two
subgroup \(\mathcal Z\), from actual Galois elements,
and this element outside it. It is therefore the
whole product. The equality for \(\ker\Delta\cap\mathcal H\)
follows immediately from its parity description.
For a general \(a\in\ker\Delta\), choose \(\sigma\)
as in (C.18m). There is \(\tau\in G_{\mathbb Q^{\rm ab}}\)
with \(\Theta(\tau)=\Theta(a\sigma^{-1})\), by (C.18d).
The full frame homomorphism then gives
\(\widetilde\Theta(a)=\widetilde\Theta(\tau\sigma)\).
The converse inclusion holds because \(\Delta\)
vanishes on Galois. This proves (C.18n).
This proves the corrected form of Lead 7 completely.
\end{proof}

\subsection{The icosahedral inference and the literature}

\begin{statement}{C.19}{The multiplier-one group is torsion-free}{Proved}
The group
\(\mathcal H=\ker(\lambda:\widehat{\rm GT}_0
\to\widehat{\mathbb Z}^{\times})\)
has no nonidentity element of finite order.
Consequently a finite subgroup
\(\mathcal K\leq\mathcal H\) is trivial.
If \(\mathcal K\) contains an element
with \(\Delta=1\), it is infinite.
\end{statement}
\begin{proof}
We prove the multiplier-one assertion
needed here; the stronger global
multiplier classification in R15 is
not needed. Suppose a nonidentity
normalized positive peripheral
automorphism \(\alpha\) of
\(P=\widehat F(x,y)\) has prime
order \(p\). A finite-order gentle
element gives one after taking a
power: its normalized action is
a genuine automorphism, and the
pair-to-action map is injective.
It acts identically on abelianization.

First its action on \(P^{(p)}\)
is identity. At the first nonidentity
lower central correction, the
\(p\)-fold iterate gives \(p\)
times that correction.
The graded pieces of a free pro-\(p\)
group are free \(\mathbb Z_p\)-modules,
by the Hall commutators in the free
nilpotent quotients, and have no
\(p\)-torsion. Induction makes every
correction vanish. These quotients
separate \(P^{(p)}\), since every
finite \(p\)-group is nilpotent.
We do not infer that \(\alpha\)
is trivial on \(P\) from this fact.
The rest of the proof supplies the
full profinite faithfulness step.

Set \(E=P\rtimes\langle\alpha\rangle\).
For \(c_i=x,y,(xy)^{-1}\), its
procyclic normalizer in \(E\)
surjects to \(C_p\), with kernel
\(\langle c_i\rangle\simeq\widehat{\mathbb Z}\).
Here
\(C_P(c_i)=\langle c_i\rangle\)
follows from the free-product
centralizer theorem
\citep[Theorem B', p.~159]{cHerfortRibes1985};
each \(c_i\) is a free basis element.
A normalizer acts trivially on
the nonzero primitive abelian class,
so it is also a centralizer.
Choose \(a_i\) above \(\alpha\).
Positive peripherality gives
\(a_ic_ia_i^{-1}=c_i\), and
\(a_i^p=c_i^{m_i}\).
The quotient
\[
 E\longrightarrow P^{(p)}\times C_p
\]
sends \(a_i\) to
\((c_i^{v_i},\alpha)\), by the
pro-\(p\) cusp-centralizer theorem.
Thus \((m_i)_p=pv_i\).
Division by \(p\) is possible in
every other component of
\(\widehat{\mathbb Z}\), so choose
\(u_i\) with \(pu_i=-m_i\).
Then \(a_ic_i^{u_i}\) has order
\(p\) and centralizes the entire
cusp group; call its subgroup \(S_i\).

Apply the profinite Brown restriction
theorem in
\citet[Theorem 2 and \S1.1,
equations (3)--(6)]{cKochloukovaZalesskii2007}.
Its hypotheses hold: \(P\) is open
in \(E\) and has \(p\)-cohomological
dimension one; finite subgroups
inject into \(E/P=C_p\), since a
free profinite group is torsion-free,
so \(E\) contains no \(C_p^2\).
Hochschild--Serre gives
\(\dim H^n(E,\mathbb F_p)\le3\)
for \(n>1\).
For each subgroup \(S\) of order
\(p\), conjugation on \(S\) is
trivial because its projection to
the abelian quotient \(C_p\) is
injective. Consequently
\(N_E(S)=C_P(S)\times S\).
Brown restriction is a dense
embedding into the product of
normalizer cohomology groups.
Its finite-dimensional source makes
it onto, and there can be only
finitely many nonzero factors.
K\"unneth and \(cd_p C_P(S)\le1\)
therefore give
\[
 \sum_{[S]}
 \bigl(1+\dim H^1(C_P(S),\mathbb F_p)\bigr)
       \le3.
 \tag{C.19a}
\]
Each peripheral \(S_i\) contributes
at least two if it is a new class:
its fixed subgroup contains a cusp
whose mod-\(p\) abelianization is
nonzero. Hence the three \(S_i\)
are conjugate. Conjugacy of these
complements can be effected by \(P\):
write \(E=P S_i\) and remove the
factor in \(S_i\) from a conjugator.
Their common fixed subgroup \(C\)
contains conjugates of all three
cusps. The first two have independent
images in \(H_1(P,\mathbb F_p)\),
so \(\dim H^1(C,\mathbb F_p)\ge2\).
Equation (C.19a) gives equality.

We now use the entire procyclic cusp
groups retained above. If the closed
subgroup \(C\) were proper, choose a
proper open subgroup \(V\) containing
it. This open subgroup corresponds to
a finite connected cover of the
three-holed sphere, of degree
\(d=[P:V]>1\). The three chosen
conjugate cusp groups give three
distinct boundary components of
degree one. To justify this for a
profinite conjugator, replace it by
a discrete element in the same
\(V\)-coset; the two resulting
elements are conjugate inside \(V\)
and have the same abelian class.
The coset fixed by the entire cusp
group is a boundary cycle of length
one in the finite covering.

Above each base boundary the sum of
the boundary degrees is \(d\).
Since \(d>1\), the three selected
components cannot be all the
boundary components. On an oriented
compact surface with boundary, their
homology classes have only the one
relation given by the sum of all
boundary classes. Therefore these
three selected classes are independent
in \(H_1(V,\mathbb F_p)\). Choose
three characters of \(V\) dual to
them. Restriction to \(C\), which
contains the three actual cusp
generators, leaves these characters
independent. This contradicts
\(\dim H^1(C,\mathbb F_p)=2\).
Hence \(C=P\).

The corresponding contribution to
(C.19a) is three, leaving a single
section class, including the
original \(\langle\alpha\rangle\).
Conjugating a section whose fixed
group is all \(P\) still has fixed
group all \(P\). Thus \(\alpha\)
is the identity on \(P\), contrary
to its assumed order.
This is a full profinite argument,
not detection by the pro-prime
quotients alone.
The subgroup conclusions follow
since \(\Delta=1\) ensures that
the chosen element is nonidentity.
\end{proof}

\begin{statement}{C.20}{The remaining icosahedral membership tests}{Open}
For \(K=\Gamma(5)/\{\pm1\}\subset B\),
membership of either selected full
lift \(e,d\) in the \(\mathcal K\)
of Lead 8 is not decided here.
C.19 proves the finiteness inference,
but does not prove these membership
tests or the claimed normalizer
identification in the attached note.
\end{statement}

The test concerns the entire
restriction to \(\widehat K\)
and its mapping-class subgroup
\(\widehat\Gamma_{0,[12]}\).
CM coordinates alone do not give
this comparison. It has not been
shown that the compactness fibres
of the full lifts meet, or avoid,
that subgroup. Membership of a
lift with \(\Delta=1\) would
force infinitude by C.19.
Nonmembership alone implies
neither finiteness of \(\mathcal K\)
nor vanishing of \(\Delta\) on it.

We checked a later possible route rather than stopping
at the attached note. The published version of
\citet[Theorem C(iii), Proposition 2.4]{cBoggi2026}
asserts that the profinite centralizer of every finite
subgroup is the closure of its discrete centralizer.
For the icosahedral \(A_5\) this assertion would settle
the question: its quotient orbifold has three points
with distinct stabilizer orders \(2,3,5\), so a
centralizing mapping class descends to the identity
on the three-point sphere. Its remaining deck
transformation belongs to \(Z(A_5)=1\). Thus its
discrete centralizer is trivial.

We state precisely why we cannot use the cited
assertion as an audited proof. Proposition 2.4's
displayed argument, in both the opened arXiv revision
and the published article, uses
\[
 \bigcap_i\overline{C_G(x_i)}
      =\overline{\bigcap_i C_G(x_i)}.
 \tag{C.20a}
\]
Only the inclusion from right to left is automatic.
The invoked result
\citep[Theorem A(iii)]{cBoggiZalesskii2025}
identifies each prime-power centralizer with its
closure, but supplies no compatibility of the
approximating conjugators for different primes.
The source gives no additional argument for (C.20a).
This is an unproved step in the proof used for this
route, not a counterexample to its conclusion.

For completeness the exact conditional consequence
does not require the note's whole normalizer
dictionary. A multiplier-one gentle element induces
an inner automorphism on \(A_5\), since its congruence
intertwiner modulo five has determinant one. Correct
by an inner element of \(\mathcal B\) with that
quotient value. Its restriction to \(\widehat K\)
then centralizes the deck \(A_5\) in
\(\operatorname{Out}(\widehat K)\). If that restriction
is a profinite mapping class and the asserted
centralizer equality holds for this \(A_5\), it is
trivial. Correcting further by an element of
\(\widehat K\) gives an automorphism identity on
\(\widehat K\), hence identity on \(\mathcal B\):
its discrepancy on any element centralizes
\(\widehat K\), whose centralizer is trivial.
Indeed that centralizer injects into the finite
quotient \(A_5\), while the centralizer of a
nonidentity torsion element of \(C_2\amalg C_3\)
is a conjugate cyclic factor and cannot contain
\(\widehat K\).

Finally the normalized gentle action is faithful
in \(\operatorname{Out}(\mathcal B)\). If it were
inner, its conjugating element would lie in \(F\)
because the action on \(S_3\) is identity and
\(Z(S_3)=1\). Fixing \(x=\tau^{-2}\) makes that
element \(x^u\). Comparing the second generator
then gives \(f=y^v x^{-u}\); its zero abelianization
forces \(u=v=0\).
Thus the extra centralizer equality would imply
\(\mathcal K=1\), exclude both \(e,d\), and make
\(\Delta|_{\mathcal K}=0\). Without a proof of
(C.20a) for the actual icosahedral action, these
remain conditional, with exactly the missing
comparison identified.

\begin{statement}{C.21}{Verified literature scope for task C.8}{Known}
Harbater--Schneps identify the
four-point symmetric peripheral
outer group with the gentle group
and the corresponding higher
genus-zero groups with full GT.
Enriquez proves full-GT cyclic-cover
compatibility; Marin constructs
the congruence cocycle and its
Kummer specialization;
Dolgushev--Guynee reconstruct the
gentle group from all finite
shadows.
\end{statement}
\begin{proof}[Derivation and source check]
We opened the primary texts
\citep[Main Theorem, \S0.1]{cHarbaterSchneps2000},
\citep[Theorem 8.1]{cEnriquez2007},
\citep[\S4]{cMarin2008}, and
\citep[Theorem 5.2, Corollary 5.4,
Proposition 5.5]{cDolgushevGuynee2024}.
The first distinguishes the
three-strand conditions from
the extra pentagon; it does
not equate the two groups.
The second supplies C.15's
actual cyclic-cover maps.
The third gives the shear
normalization in C.3 and C.16.
The full author copy of
\citet[Propositions 4.4, 4.13 and Theorem 4.16]{cNakamuraTsunogaiYasuda2010}
was also opened. Its two scopes differ:
the triangle-cover beta action is proved on full GT,
while the modular equianharmonic-frame comparison
is proved on the Galois image. The introduction
explicitly leaves the broader power-conjugacy question
undecided there. C.15 records the resulting exact
dependency rather than extending those quantifiers.
For the fourth, a genuine
finite shadow lies in the
image from every finer level.
Conversely the finite fibres
have the finite intersection
property, giving a compatible
inverse-limit pair. Generation
at every finite quotient gives
invertibility. This explains
why one finite CM shadow is
insufficient. LM7 D.25's
twelve selected extra shadows
are removed by LM7 D.27;
they are not E.58's full lifts.

Finite characters exist already
by composing the multiplier with
finite characters of
\(\widehat{\mathbb Z}^{\times}\).
Marin's additive character on
the multiplier kernel also
gives finite characters after
reduction. Neither observation
classifies all finite characters
or resolves R14's particular
nonarithmetic character.
The free-product centralizers
of \citet{cHerfortRibes1985}
and Brown restriction of
\citet{cKochloukovaZalesskii2007}
have the precise uses in C.19;
they do not classify torsion
conjugacy in either GT group.

We also correct a source
reference in the attached
LM4 D.32 referee report.
In the primary author version
of \citet{cLochakSchneps2006},
question 1.4 asks whether
complex conjugation is
self-normalizing and records
an affirmative solution.
It is not the torsion-conjugacy
question attributed to that
number in the report.
A different intended version
must be identified explicitly.
\end{proof}

\begin{statement}{C.22}{Limit of the literature classification search}{Open}
This return does not establish a
complete classification of finite
characters or torsion conjugacy
classes in the two profinite GT
groups, or an unconditional
literature theorem distinguishing
them. No priority is claimed.
\end{statement}

On 16 November 2026 we searched
both available engines for
profinite strict containment,
gentle versus full GT, finite
characters, and finite-order
elements, and opened the
primary texts listed above.
We also inspected contemporary
gentle-shadow work and the
supplied referee reconstructions.
For the icosahedral problem we also opened
\citet{cBoggiZalesskii2025} and the published
September 2026 version of \citet{cBoggi2026};
C.20 records the exact intersection step still
needed by that apparent solution.
These retrieved sources did
not settle the outstanding
global image and CM-parity
questions. This is a bounded
search report, not a proof
that no later theorem exists.
Completed profinite groups
remain distinct from their
rational pro-unipotent
comparison groups throughout.
